\section{技术转型与 Wedge：为什么是 Cloud Identity}

McKinnon 把创业机会放在 cloud transition 中理解：企业应用从自建机房转向 SaaS 后，旧网络边界和本地目录不再自然覆盖所有资源。新公司无法在 incumbent 最擅长的旧市场正面复制，而应找到转型暴露的新 control point，再用一个客户愿意立刻付费的 wedge 切入。

\subsection{从 Monitoring Pivot 到 SSO}

团队先做 cloud application monitoring，但 buyer feedback 平淡；相反，员工面对 Gmail、Dropbox 等多个 SaaS 账户时，重复登录、密码和离职权限成为即时痛点。\term{single sign-on}（单点登录，SSO）让用户通过一个受信任 IdP 访问多个应用，减少重复凭据并集中 policy。SSO 只是 wedge，却天然要求与许多应用、目录和协议连接。

\lecturefigure{01-transition-window.png}{技术转型会制造新的控制点；identity wedge 从即时登录痛点成长为共享基础设施。}{本地字幕 00:00--05:20；概念重绘。}

\begin{knowledgebox}{读图：Transition、pain 与 platform 要同时成立}
Cloud transition 改变应用位置与购买方式；credential fragmentation 形成具体 pain；SSO 提供短期可付费结果；directory、policy 和 lifecycle 则让产品拥有平台深度。只有 transition 没有 buyer pain，会变成长期研究；只有小痛点没有扩展面，会成为 feature；只有 platform vision 没有 wedge，则很难获得第一批生产流量。
\end{knowledgebox}

\teachervoice{McKinnon 没有把第一个想法浪漫化：monitoring prototype 做出来后，潜在客户并不兴奋，团队因此转向“让 SaaS 登录更简单”。课堂提示是先寻找真实 willingness to pay，再证明小入口后面存在足够深的 infrastructure。}

\subsection{基础设施冰山：简单体验背后的复杂性}

一个员工点击应用图标，看起来只是“少输一次密码”；系统却必须同步用户状态、选择身份源、执行认证策略、签发 assertion/token、建立 session、处理 logout 和 revoke，并在组织变更时更新 entitlements。Identity wedge 的复利来自这些共享 primitive 可服务越来越多应用，而不是来自登录页面本身。

\begin{importantbox}{Wedge 的两个约束}
一个基础设施创业入口同时需要：\textbf{today value}，客户今天就愿意为更少登录摩擦、更快 onboarding 或更可控 offboarding 付费；\textbf{tomorrow depth}，入口能够扩展到 directory、MFA、policy、provisioning、governance 与 threat response。两者缺一，产品要么没有现金流，要么没有长期平台空间。
\end{importantbox}

\subsection{本章小结}

Cloud identity 的机会来自应用边界变化。SSO 是可卖的前门，identity lifecycle 与 policy 才是后续平台。Pivot 的依据是客户证据，而不是创始人对第一个 prototype 的沉没成本。

